% Four new mixed van der Waerden numbers
% LaTeX source; compiles with pdflatex, xelatex, or tectonic.
\documentclass[11pt]{article}

\usepackage[margin=1in]{geometry}
\usepackage{amsmath,amssymb,amsthm}
\usepackage{booktabs}
\usepackage{microtype}
\usepackage[hidelinks]{hyperref}

\newtheorem{theorem}{Theorem}
\theoremstyle{remark}
\newtheorem{remark}{Remark}

\newcommand{\seq}[1]{\texttt{#1}}

\title{Four new mixed van der Waerden numbers}
\author{Leo Y. Zhang}
\date{August 2026}

\begin{document}

\maketitle

\begin{abstract}
We establish four previously unknown mixed van der Waerden numbers:
$w(14;\,2^{12},3,4)=57$, $w(21;\,2^{19},3,3)=52$, $w(9;\,2^{7},4,4)=68$ and
$w(6;\,2^{4},4,5)=84$. Each extends by one term a sequence in the On-Line
Encyclopedia of Integer Sequences (\seq{A217058}, \seq{A217005},
\seq{A217007}, \seq{A217236}) whose published values, due to Ahmed, had stood
since 2012. Every lower bound ships as an explicit colouring, checkable
against the definition in milliseconds by a standalone program that uses only
the Python standard library and never invokes a SAT solver. Every upper bound
is a SAT refutation obtained by cube-and-conquer and re-derived along
independent paths chosen so that they cannot share a mistake, including one
engine that applies no symmetry breaking at all. For the first of the four
values the refutation has additionally been reduced to formally checked proof
objects: $4{,}487$ per-cube DRAT proofs, each replayed to \texttt{s VERIFIED}
by \texttt{drat-trim} against a single formula, together with a checked proof
that the cube set is exhaustive. The refutations behind the remaining three
values are cross-checked but have not been reduced to proof objects. A fifth
value, $w(11;\,2^{9},3,5)=74$, was computed but is deliberately withheld
because one of our own validation gates for its family did not run to
completion; we state this as policy. All code, certificates, proof-object
tooling and result records are public.
\end{abstract}

\noindent\emph{2020 Mathematics Subject Classification:} 05D10; 68R05.

\section{Introduction}\label{sec:intro}

Van der Waerden's theorem \cite{vdw1927} states that for any positive
integers $t_1,\dots,t_r$ there is a least $n$, the van der Waerden number
$w(r;t_1,\dots,t_r)$, such that every $r$-colouring of $\{1,\dots,n\}$
contains, for some $i$, a $t_i$-term arithmetic progression entirely in
colour $i$. Exact values are notoriously scarce; see Landman and Robertson
\cite{landmanrobertson} for background. The classical diagonal case
$w(2;6,6)=1132$ was settled by Kouril and Paul \cite{kourilpaul}, and
SAT-based methods have been the dominant tool for exact van der Waerden
computations since Dransfield et al.\ \cite{dransfield} demonstrated their
reach; see also Ahmed, Kullmann and Snevily \cite{aks} for the off-diagonal
family $w(2;3,t)$.

This paper concerns the \emph{mixed} numbers
$w(j+r;\,2^j,t_1,\dots,t_r)$, in which $j$ of the colour classes have target
$2$. A class with target $2$ can hold at most one element, since any two
integers form a $2$-term progression, so these classes behave as a budget of
$j$ \emph{wildcards}: positions excused from the problem at a cost of one
budget unit each. Ahmed computed extensive tables of such numbers
\cite{ahmed2009,ahmed2011,ahmed2013} and entered several families into the On-Line
Encyclopedia of Integer Sequences in 2012 \cite{oeis}. The four families
treated here --- \seq{A217058} $=w(j+2;\,2^j,3,4)$, \seq{A217005}
$=w(j+2;\,2^j,3,3)$, \seq{A217007} $=w(j+2;\,2^j,4,4)$ and \seq{A217236}
$=w(j+2;\,2^j,4,5)$ --- had received no new terms since then.

\subsection{Results}

We establish the next term of each of the four families.

\begin{center}
\begin{tabular}{llcl}
\toprule
sequence & new term & value & status \\
\midrule
\seq{A217058} & $a(12) = w(14;\,2^{12},3,4)$ & $57$ & approved by the OEIS, August 2026 \\
\seq{A217005} & $a(19) = w(21;\,2^{19},3,3)$ & $52$ & submitted, under review \\
\seq{A217007} & $a(7) \;\,= w(9;\,2^{7},4,4)$ & $68$ & submitted, under review \\
\seq{A217236} & $a(4) \;\,= w(6;\,2^{4},4,5)$ & $84$ & submitted, under review \\
\bottomrule
\end{tabular}
\end{center}

A fifth value, $w(11;\,2^9,3,5)=74$ (the next term of \seq{A217059}), was
computed by the same method and passes the same certificate checks, but is
\emph{withheld} rather than claimed: the validation gate that reproduces the
published adjacent term of that family as an independent check on the method
was started and killed without a verdict. Section~\ref{sec:withheld} states
the policy.

\subsection{The certificate discipline}

The two halves of each value have very different epistemic status, and the
entire design of this computation follows from taking that seriously.

A \emph{lower} bound asserts the existence of a finite object. We ship the
object: an explicit colouring, printed in full in this paper, which anyone
can check against the definition by hand or with a short standalone verifier
(\texttt{vdw/verify\_certificate.py} in the repository) that uses only the
Python standard library and shares no code with the search. The lower bounds
require no trust in anything reported here.

An \emph{upper} bound asserts the absence of an object, and a SAT solver's
\textsc{unsat} answer is only as good as the claim that the formula given to
the solver faithfully represents the problem, plus the solver itself. We
address the first risk --- the far likelier one --- by proving, on every
instance small enough to enumerate exhaustively, that the CNF accepts exactly
the colourings the definition accepts, in both directions
(Section~\ref{sec:verification}); and the second by re-deriving every
refutation along independent paths, and, for the headline value, by reducing
the refutation to DRAT proof objects checked by an independently built
verifier (Section~\ref{sec:drat}).

We are explicit throughout about which claims carry which kind of evidence.
In particular: of the four new values, only the refutation behind
$w(14;\,2^{12},3,4)=57$ has been reduced to formally checked proof objects;
the other three rest on cross-checked solver verdicts, not on checked proofs.

\section{Definitions and conventions}\label{sec:defs}

Let $t_1,\dots,t_r \ge 2$ be integers. A colouring of $\{1,\dots,n\}$
assigns each integer to one of $r$ classes. It is \emph{valid} for
$(t_1,\dots,t_r)$ if class $i$ contains no $t_i$-term arithmetic progression,
for every $i$. The van der Waerden number $w(r;t_1,\dots,t_r)$ is the least
$n$ for which no valid colouring of $\{1,\dots,n\}$ exists; van der Waerden's
theorem guarantees it is finite.

Writing $2^j$ for $j$ classes of target $2$, the mixed numbers of interest
are
\[
  w(j+r;\,2^j,t_1,\dots,t_r).
\]
Since a class with target $2$ holds at most one element, an equivalent
formulation is
\[
  w(j+r;\,2^j,t_1,\dots,t_r) = 1 + \max\bigl\{ n : \{1,\dots,n\}
  \text{ admits a valid colouring with at most } j \text{ wildcards} \bigr\},
\]
where a colouring with wildcards assigns each integer either to one of the
$r$ real classes or to a wildcard position that belongs to no class. In
certificates below, \texttt{.}\ denotes a wildcard and \texttt{1},
\texttt{2} the two real classes.

For a family with fixed targets we write $a(j) = w(j+2;\,2^j,t_1,t_2)$,
following the OEIS entries, which index each family by the number of
wildcards with offset $0$.

\section{Method}\label{sec:method}

\subsection{Encoding}\label{sec:encoding}

For given $n$ and $j$, introduce Boolean variables $v(i,c)$ for
$i \in \{1,\dots,n\}$ and $c \in \{0,1,\dots,r\}$, where $c=0$ means
``wildcard'':
\begin{itemize}
\item exactly one class per position: one clause $\bigvee_c v(i,c)$ and
  pairwise exclusions;
\item no monochromatic progression: for each class $c$ with target $t_c$ and
  each $t_c$-term arithmetic progression $A \subseteq \{1,\dots,n\}$, the
  clause $\bigvee_{i \in A} \neg v(i,c)$;
\item wildcard budget: a totalizer encoding \cite{totalizer} of
  $\sum_i v(i,0) \le j$.
\end{itemize}
The formula is satisfiable exactly when a valid colouring exists, so $a(j)$
is the least $n$ whose formula is unsatisfiable. Solving is done through
PySAT \cite{pysat} with CaDiCaL~1.9.5 \cite{cadical} as the CDCL back end;
proof-object generation (Section~\ref{sec:drat}) uses a standalone
Kissat~4.0.1 \cite{cadical}, for reasons given there.

\subsection{Monotonicity and a free lower bound}\label{sec:mono}

If $\{1,\dots,n\}$ admits a valid colouring, so does $\{1,\dots,m\}$ for
every $m<n$: restrict the colouring; no progression is created by deleting
elements, and the wildcard count cannot increase. Satisfiability is therefore
monotone decreasing in $n$ and $a(j)$ is the unique threshold. This licenses
probing candidate values directly --- any \textsc{sat} answer raises the
floor, any \textsc{unsat} answer caps the value --- rather than climbing one
$n$ at a time, and it lets the two halves of a value run as independent
concurrent computations.

Monotonicity in $j$ gives a bound for nothing: $a(j+1) \ge a(j)+1$. A valid
colouring of $\{1,\dots,a(j)-1\}$ with at most $j$ wildcards extends to
$\{1,\dots,a(j)\}$ by making the new final position a wildcard; it uses at
most $j+1$ wildcards and creates no monochromatic progression. Whether this
free bound is tight varies by family, and we report its role honestly in
each case: for \seq{A217007} it \emph{is} the lower bound
(Section~\ref{sec:a217007}); for \seq{A217059} it was three short and for
\seq{A217236} four short, so those witnesses are genuine search results.

\subsection{Reversal symmetry}\label{sec:revsym}

The map $i \mapsto n+1-i$ is an automorphism of the problem: it carries the
$t$-term progression $a, a+d, \dots, a+(t-1)d$ to
$n+1-a-(t-1)d, \dots, n+1-a$, again a progression with the same common
difference; it fixes every class and the wildcard count. Requiring a
colouring to be lexicographically no greater than its own reversal is
therefore a sound lex-leader constraint, and it halves the search space.

This matters more than it might appear. The pre-existing implementation broke
only the symmetry between classes \emph{sharing} a target value. On
distinct-target families such as $(3,4)$ that rule emits no clauses at all,
so every search there had previously run with no symmetry breaking
whatsoever. Adding reversal symmetry measured a $1.55\times$ speedup on the
$n=45$, $j=8$ refutation, on the half of the problem that consumes
essentially all the time. Because this constraint is the one piece of new
mathematics in the engine, the verification programme of
Sections~\ref{sec:verification} and~\ref{sec:drat} is arranged so that it is
never assumed by its own check.

\subsection{Cube-and-conquer and the crash-honest rule}\label{sec:cube}

Refutations use cube-and-conquer \cite{cubeandconquer}: branch on all class
assignments to the first $k$ positions, discard prefixes that already exceed
the wildcard budget or already contain a monochromatic progression, and solve
the residual formulas in parallel. At $k=4$ the four families yield $75$
cubes for targets $(3,4)$, $36$ for $(3,3)$, $40$ for $(4,4)$ and $80$ for
$(4,5)$. A measured comparison on $n=45$, $j=8$ put $k=4$ at $34.1$\,s
against $k=6$ at $65.1$\,s, so the finer split was not used for search.

One structural rule is worth stating because it is the difference between a
theorem and a retraction: \textsc{unsat} is reported only when every cube has
returned an explicit verdict. A worker killed by the operating system raises
an error rather than being counted as an empty branch. Three earlier runs
were lost to precisely that class of failure.

\section{The four new terms}\label{sec:results}

\subsection{$\seq{A217058}$: $w(14;\,2^{12},3,4)=57$}\label{sec:a217058}

One class must avoid $3$-term progressions, the other $4$-term progressions.
The published terms, due to Ahmed \cite{ahmed2009,ahmed2013}, are
\[
a(0..11) = 18,\,21,\,25,\,29,\,33,\,36,\,40,\,42,\,45,\,48,\,52,\,55 ,
\]
ending at $w(13;\,2^{11},3,4)=55$. The OEIS entry was queried directly on
2026-07-28 and returned exactly these twelve terms with an auto-synthesised
b-file, i.e.\ no extension had been contributed.

\begin{theorem}\label{thm:a217058}
$a(12) = w(14;\,2^{12},3,4) = 57$.
\end{theorem}

\emph{Lower bound.} The following colouring of $\{1,\dots,56\}$ uses exactly
$12$ wildcards, has no $3$-term progression in class~\texttt{1} and no
$4$-term progression in class~\texttt{2}, so $a(12) > 56$:

\begin{center}\small
\texttt{2.21221212.12.22211.112.2221.222.2.1..12221211212..22212}
\end{center}

\noindent Class~\texttt{1} occupies $16$ positions, class~\texttt{2} occupies
$28$, and $12$ are wildcards. The standalone verifier checks it in
milliseconds.

\emph{Upper bound.} No valid colouring of $\{1,\dots,57\}$ with at most $12$
wildcards exists. The refutation was established along paths chosen so that
they cannot share a mistake:

\begin{center}
\begin{tabular}{lllclr}
\toprule
encoding & symmetry & solver & cubes & verdict & time \\
\midrule
vdw4 & reversal on & CaDiCaL 1.9.5 & $k=4$ & \textsc{unsat} & $6257.1$\,s \\
vdw2 & none at all & CaDiCaL 1.9.5 & $k=4$ & \textsc{unsat} & $8036.6$\,s \\
vdw4 & reversal off & CaDiCaL 3.0.0 & $k=5$ & (not completed) & --- \\
vdw4 & reversal off & CaDiCaL 1.9.5 & $k=4$ & \textsc{unsat} & $8317.4$\,s \\
\bottomrule
\end{tabular}
\end{center}

The second row is the one that matters: \texttt{vdw2}, an older and
independently written engine, imposes no lex-leader constraint on this family
and searches the full unreduced space, so it cannot inherit an error from the
reversal-symmetry constraint. It agrees. The third row was started and
stopped after $7.4$ hours without finishing; it is reported as incomplete
rather than omitted. The same reversal-off question was settled in the
configuration that suits this family (row four), and a second \texttt{vdw2}
refutation on that pass agrees in $7259.8$\,s. Separately, a
randomised-restart portfolio --- $8$ rounds $\times$ $5$ seeds $\times$
$3{,}000{,}000$ conflicts, with initial phases drawn from the class
distribution of real witnesses --- hunted for a witness at $n=57$ and found
none. Section~\ref{sec:drat} reduces this refutation to checked proof
objects.

\begin{remark}
$a(12)-a(11)=2$, while the published differences are
$3,4,4,4,3,4,2,3,3,4,3$, so a naive extrapolation predicts $58$ or $59$;
both were refuted before $57$ was reached. A step of $2$ already occurs at
$a(6)=40 \to a(7)=42$. The disagreement with extrapolation was treated as
grounds for additional verification rather than as a curiosity.
\end{remark}

\subsection{$\seq{A217005}$: $w(21;\,2^{19},3,3)=52$}\label{sec:a217005}

Both classes must avoid $3$-term progressions. Published terms (offset $0$,
last extended by Ahmed in December 2012, confirmed live against the OEIS API
on 2026-07-29):
\[
a(0..18) = 9,\,14,\,17,\,20,\,21,\,24,\,25,\,28,\,31,\,33,\,35,\,37,\,39,\,42,\,44,\,46,\,48,\,50,\,51 .
\]

\begin{theorem}\label{thm:a217005}
$a(19) = w(21;\,2^{19},3,3) = 52$.
\end{theorem}

\emph{Lower bound.} This colouring of $\{1,\dots,51\}$ uses exactly $19$
wildcards and contains no $3$-term progression in either class, so
$a(19) > 51$:

\begin{center}\small
\texttt{..11.1122.2211.1122.22.........11.1122.2211.1122.22}
\end{center}

\noindent It was found by search in $1520.1$\,s and accepted by the
standalone verifier. Note also that the free bound of
Section~\ref{sec:mono} gives $a(19) \ge 52$ independently of any solver, so
the refutation at $n=52$ is the only computational input to the value.

\emph{Upper bound.} $n=52$, $j=19$ is unsatisfiable: $4507.0$\,s, all $36$
cubes reporting an explicit verdict. Cross-checks: \texttt{vdw2} refuted the
same instance in $10650.4$\,s, the primary engine with reversal symmetry
switched off in $11779.0$\,s, and \texttt{vdw2} additionally found and
verified a witness of its own at $n=51$ ($1934.0$\,s).

The targets here are equal, which switches on the colour-swap symmetry
breaker --- a code path the first family never touches. On equal-target
families the two breakers together are sound but \emph{not} canonical: they
retain two representatives per orbit rather than one, because composing two
independent lex-leader constraints does not canonicalise the group they
generate. That costs time and cannot cost correctness, and the audit of
Section~\ref{sec:verification} confirms zero orbits are lost. The family's
full-scale gate was run before the extension was attempted: published
$a(18)=51$ was reproduced exactly --- \textsc{sat} at $n=50$ with a verified
witness and \textsc{unsat} at $n=51$, both at $j=18$, in $3469.4$\,s.

\begin{remark}
The witness is strikingly structured: the segment
\texttt{11.1122.2211.1122.22} appears twice, separated by a block of nine
consecutive wildcards. Whether the extremal colourings of this family are
genuinely periodic is not something a single example can settle, and no claim
is made here; it is recorded as the sort of structure worth examining if
anyone pursues the family further.
\end{remark}

\subsection{$\seq{A217007}$: $w(9;\,2^{7},4,4)=68$}\label{sec:a217007}

Both classes must avoid $4$-term progressions. Published terms (offset $0$,
OEIS keyword \texttt{hard}, confirmed live against the OEIS API on
2026-07-29):
\[
a(0..6) = 35,\,40,\,53,\,54,\,56,\,66,\,67 .
\]

\begin{theorem}\label{thm:a217007}
$a(7) = w(9;\,2^{7},4,4) = 68$.
\end{theorem}

\emph{Lower bound.} This colouring of $\{1,\dots,67\}$ uses exactly $7$
wildcards and contains no $4$-term progression in either class, so
$a(7) > 67$:

\begin{center}\small
\texttt{..1112112111.2221221222.1112112111.2221221222.1112112111.2221221222}
\end{center}

An honest qualification: this certificate \emph{is} the free construction of
Section~\ref{sec:mono} --- the family gate's $a(6)$ colouring with one
further wildcard prepended. The solver, given $n=67$, $j=7$ cold,
rediscovered that construction rather than finding a structurally different
witness. It is reported as a verified confirmation, not as independent
evidence, and the entire computational weight of the theorem sits on the
refutation at $n=68$.

\emph{Upper bound.} $n=68$, $j=7$ is unsatisfiable: $7269.0$\,s, all $40$
cubes reporting an explicit verdict. Cross-checks: \texttt{vdw2}, which on
this equal-target family breaks only the colour swap and carries no
reversal-symmetry constraint, refuted the instance in $7375.6$\,s, and the
primary engine with reversal symmetry switched off refuted it in
$7285.2$\,s. The family gate passed first: published $a(6)=67$ was reproduced
exactly --- \textsc{sat} at $n=66$ with a verified witness using all $6$
wildcards and \textsc{unsat} at $n=67$, both at $j=6$, $1032.1$\,s for the
pair.

\begin{remark}
The step $a(6)=67 \to a(7)=68$ is $+1$, the smallest possible; $+1$ already
occurs twice in the published sequence, at $a(2)=53 \to a(3)=54$ and at
$a(5)=66 \to a(6)=67$. The witness is conspicuously periodic --- the block
\texttt{1112112111.2221221222} repeats three times, and the two classes are
exact complements under $\texttt{1} \leftrightarrow \texttt{2}$ --- but
since the certificate is inherited from the $a(6)$ gate rather than found
independently, this says something about that colouring, not about extremal
colourings of the family in general.
\end{remark}

\subsection{$\seq{A217236}$: $w(6;\,2^{4},4,5)=84$}\label{sec:a217236}

Class~\texttt{1} must avoid $4$-term progressions, class~\texttt{2} must
avoid $5$-term progressions. Published terms (offset $0$, OEIS keyword
\texttt{hard}, contributed by Ahmed in September 2012 and confirmed live
against the OEIS API on 2026-07-29):
\[
a(0..3) = 55,\,71,\,75,\,79 .
\]
Four terms is the shortest published list of the families treated here and
$j=4$ the smallest wildcard budget, but $n$ is the largest: the refuted
instance is $\{1,\dots,84\}$.

\begin{theorem}\label{thm:a217236}
$a(4) = w(6;\,2^{4},4,5) = 84$.
\end{theorem}

\emph{Lower bound.} This colouring of $\{1,\dots,83\}$ uses exactly $4$
wildcards, with no $4$-term progression in class~\texttt{1} and no $5$-term
progression in class~\texttt{2}, so $a(4) > 83$:

\begin{center}\small
\texttt{122121221221212221.212121221121222211121.221212222.2222.212211211122221211122212122}
\end{center}

\noindent Unlike the previous family's, this lower bound was earned rather
than inherited: the free construction from $a(3)=79$ gives only
$a(4) \ge 80$, four short of the truth. The witness above was found by
search ($6469.7$\,s, cube $52$ of $80$) and is an independent object ---
witnesses at $n=80$, $81$, $82$ were found in an earlier pass, and the
$n=83$ colouring extends none of them, differing from all three already at
the second position.

\emph{Upper bound.} $n=84$, $j=4$ is unsatisfiable: $7965.0$\,s, all $80$
cubes reporting an explicit verdict. The value was bracketed rather than
climbed: a ceiling probe at $n=87$ refuted in $6688.4$\,s, capping the value;
$n=83$ \textsc{sat} raised the floor to $84$; $n=85$ \textsc{unsat}
($9061.5$\,s) dropped the ceiling to $85$; and $n=84$ \textsc{unsat} fixed
the term. Monotonicity (Section~\ref{sec:mono}) is what licenses this. The
refutation was established three ways:

\begin{center}
\begin{tabular}{lllclr}
\toprule
encoding & symmetry & solver & cubes & verdict & time \\
\midrule
vdw4 & reversal on & CaDiCaL 1.9.5 & $k=4$ & \textsc{unsat} & $7965.0$\,s \\
vdw2 & none at all & CaDiCaL 1.9.5 & $k=4$ & \textsc{unsat} & $10959.3$\,s \\
vdw4 & reversal off & CaDiCaL 1.9.5 & $k=4$ & \textsc{unsat} & $16975.3$\,s \\
\bottomrule
\end{tabular}
\end{center}

\noindent \texttt{vdw2} also re-ran the satisfiable side at $n=83$ from
scratch and returned a verified witness of its own in $9892.7$\,s; it arrived
at the same colouring as above. The family gate passed before the extension
was attempted: published $a(3)=79$ reproduced exactly, \textsc{sat} at
$n=78$ ($404.9$\,s) and \textsc{unsat} at $n=79$ ($377.7$\,s), both at $j=3$.

The targets $4$ and $5$ are distinct, so as in Section~\ref{sec:a217058} the
colour-swap breaker emits no clauses and reversal symmetry is the only
breaking in force --- precisely the constraint the second and third rows
above isolate. When this term was first written up, the exhaustive encoding
audit did \emph{not} cover the target pair $(4,5)$; rather than rest on the
argument that the clause generator is generic in the targets, the audit was
extended to $(4,5)$ on 2026-07-31 and passes: the CNF accepts exactly the
colourings the definition accepts ($538$ and $7631$ on the two audit
instances, equal in both directions) and symmetry breaking loses no orbit.

\begin{remark}
The step $a(3)=79 \to a(4)=84$ is $+5$ against published differences
$16,4,4$, so an extrapolation from the last two steps would have said $83$
--- which the witness at $n=83$ refutes outright.
\end{remark}

\section{Verification of the encoding and engine}\label{sec:verification}

The upper bounds are only as strong as the claim that the formula given to
the solver represents the problem. Five independent guards were applied; all
tooling named here is in the public repository.

\subsection{The encoding equals the definition}

On instances small enough to enumerate exhaustively, the set of colourings
satisfying the CNF was compared against the set accepted by a direct
transcription of the definition that never inspects a clause. They agreed
exactly, in both directions, on all $11$ cases tested, spanning target shapes
$(3,4)$, $(3,3)$, $(3,5)$, $(3,3,3)$, $(4,4)$ and $(4,5)$ --- one shape for
every family claimed in this paper. Nothing invented, nothing lost.
(\texttt{vdw/encoding\_audit.py})

\subsection{Symmetry breaking loses nothing}

The same audit checks that every orbit of the symmetry group --- generated by
the reversal and by swaps of equal-target classes --- retains at least one
representative. Losing an orbit is exactly how a satisfiable instance would
be reported unsatisfiable. Zero orbits were lost in any case tested. On
distinct-target families, including $(3,4)$, the breaking is additionally
exactly canonical: $11{,}539$ orbits, $11{,}539$ survivors.

\subsection{The engine reproduces what is already known}

All twelve published terms of \seq{A217058} were re-derived by the same
engine that produced the thirteenth: for each $a(j)=w$ this means finding a
verified witness at $n=w-1$ \emph{and} refuting $n=w$. A further $29$
published values across five families were reproduced with the
reversal-symmetry constraint enabled. Each family's full-scale gate at the
adjacent index was run before its extension was attempted
(Sections~\ref{sec:a217005}--\ref{sec:a217236}); for \seq{A217058} the gate
is $a(11)=55$ at $j=11$: \textsc{sat} at $n=54$ with a verified witness
using all $11$ wildcards ($1872.5$\,s) and \textsc{unsat} at $n=55$
($1239.2$\,s). This gate matters because every other replayed value tops out
at $j=10$; a defect confined to large $j$ would have escaped the rest of the
testing. It did not exist.

\subsection{The cardinality constraint at full scale}

The exhaustive audit reaches only $n \le 12$ and $j \le 3$, while the
headline instance runs at $j=12$, and the totalizer's structure grows with
the bound. Too strong a cardinality constraint would forbid legal colourings
and produce precisely the kind of false refutation that matters. Tested
directly at $n=55,\dots,58$ with $j=12$, from both sides: $240$ assignments
at the limit were accepted and $240$ over the limit were rejected, and
forcing $13$ wildcards inside the full formula is correctly unsatisfiable.
(\texttt{vdw/scale\_test.py})

\subsection{Independent refutation paths}

Every established value's refutation was re-derived along disjoint paths ---
in particular through \texttt{vdw2}, an older, independently written engine
that carries no reversal-symmetry constraint --- with agreement in every
case, as detailed per family above. The records are the
\texttt{crosscheck\_*.json} files in the repository.

\section{Formally checked proof objects for the headline
refutation}\label{sec:drat}

The guards of Section~\ref{sec:verification} target an encoding defect,
which is by far the likelier failure. A defect in the solver itself is a
different risk, and the standard remedy is a DRAT proof of unsatisfiability
\cite{drattrim}: the solver logs its reasoning, and an independent checker
replays the log against the formula. What remains to be trusted shrinks to
the checker plus the encoding audit.

\subsection{Scope: what is and is not certified}

We state the perimeter plainly. The refutations behind the published rungs
$a(0)$ through $a(6)$ of \seq{A217058}, and the first rungs of the other
three families, are replayed under \texttt{drat-trim} on every run of the
repository's verification gate and come back \texttt{s VERIFIED}. The
headline refutation $n=57$, $j=12$ is certified by the cube-level route
described below. The refutations behind the other three new values ---
$w(21;\,2^{19},3,3)=52$, $w(9;\,2^{7},4,4)=68$ and $w(6;\,2^{4},4,5)=84$ ---
have \emph{not} been reduced to proof objects; those upper bounds rest on
the cross-checked solver verdicts of
Sections~\ref{sec:a217005}--\ref{sec:a217236}, not on checked proofs.

In all certified refutations, symmetry breaking is switched off: a proof of
a symmetry-broken formula does not prove the original statement, and the
reversal lex-leader constraint --- the one piece of new mathematics in the
engine --- must not be assumed by the thing meant to check it. What is
certified is the raw encoding. Proofs are generated by a standalone
Kissat~4.0.1, because the PySAT interface returns truncated proofs from
every bundled solver (the refutation never closes), and checked by
\texttt{drat-trim} \cite{drattrim} built from byte-identical upstream
source. Only the checker's literal \texttt{s VERIFIED} line counts as a
pass; every other outcome, including a parse failure, fails the gate.

\subsection{Certifying $n=57$, $j=12$ by cubes}

Proof size grows about $2.4\times$ per rung in this family, which puts a
single-shot proof of the headline instance at order $6$--$17$ hours and
$30$--$150$\,GB. Instead the search space was split at depth $k=8$ into
$4{,}487$ cubes. Each cube's refutation was produced with symmetry breaking
off and replayed to \texttt{s VERIFIED} under \texttt{drat-trim}, and all
$4{,}487$ records carry the same formula SHA-256, so they refute one formula
rather than several. The largest individual proof was under $200$\,MB.

Exhaustiveness of the cube set is proved rather than assumed
(\texttt{vdw/cube\_exhaustive.py}). The checker does not trust the cube
generator: it walks the full ternary prefix tree, and for every prefix
\emph{not} in the cube set it exhibits the clause of the formula that
refutes it --- for targets $(3,4)$ all $662$ dropped prefixes are refuted by
a monochromatic-progression clause already present, and no prefix is dropped
by the colour-symmetry rule, which for distinct targets emits nothing at
all. The resulting $3{,}236$ tail lemmas are themselves replayed by
\texttt{drat-trim} and come back \texttt{s VERIFIED}. Hence $F$ and
$F \cup \{\neg C : C \in \text{cubes}\}$ have identical models, every cube
is refuted, and the empty clause follows. This is the same compositional
pattern used at much larger scale for the Boolean Pythagorean triples
problem \cite{pythagorean}.

The pipeline is fail-closed by construction: a prefix wrongly dropped yields
a tail lemma that is not RUP, which \texttt{drat-trim} rejects; a missing
cube clause leaves an internal lemma underivable. Both failure modes were
observed failing before any positive result was believed, alongside a proof
truncated to half, a proof of a different formula, a real proof checked
against a satisfiable formula, and, for each of several published rungs, the
instance one below the published term --- which is satisfiable and must
report \textsc{sat}, never \texttt{VERIFIED}, and does.

\section{The withheld fifth value}\label{sec:withheld}

The same method produced $w(11;\,2^9,3,5) = 74$, the next term of
\seq{A217059}: a verified witness at $n=73$ found by search (the free bound
was three short, and explicit witnesses were found and verified at $n=71$,
$72$, $73$ in turn), a refutation at $n=74$ with all $76$ cubes reporting,
and cross-checks along the same disjoint paths as the other families, which
agree.

We nevertheless \emph{withhold} this value rather than claim it. Our
discipline requires every family's validation gate --- reproducing the
published adjacent term, here $a(8)=70$, as an independent check on the
method in that family --- to run to completion before an extension is
claimed, and for this family the gate was started and killed without a
verdict. Every other family in this paper passed its gate before its
extension was claimed. The computed value and its artifacts are on disk in
the public repository, marked withheld; the value has not been submitted to
the OEIS and should not be quoted as established until the gate completes.
We state this as policy: a verification step that was defined and then not
finished does not get waived retroactively because the answer looks right.

\section{Limitations}\label{sec:limitations}

These are finite, checkable extensions of tracked sequences by one term
each. They are not structural results: they say nothing about the growth of
the families and provide no new proof technique. The exhaustive
encoding-equals-definition audit reaches only instances small enough to
enumerate ($n \le 12$, $j \le 3$); beyond that scale its guarantee is
extended by the scale tests of Section~\ref{sec:verification} and by replay
of known values, not by exhaustion. Three of the four upper bounds are not
reduced to checked proof objects. All computations ran on a single 8-core
desktop machine; per-instance solver times are reported so that the scale of
the computation is not overstated.

\section{Methods and authorship}\label{sec:authorship}

The search and verification tooling, the computations, and the drafting of
this paper were carried out with substantial AI assistance, directed and
reviewed by the author. This is stated plainly because the reader is
entitled to weigh it. The verification programme of
Sections~\ref{sec:verification} and~\ref{sec:drat} is designed precisely so
that no trust in any tool, human, or model is required: every lower bound is
a printed colouring checkable against the definition by hand or by a
standard-library script that shares no code with the search; every certified
upper bound is a DRAT proof replayed by an independently built checker whose
own failure modes were tested first; and the remaining upper bounds are
labelled as what they are --- cross-checked solver verdicts. Nothing in this
paper asks to be believed on the strength of the process that produced it.

\section{Reproducibility}\label{sec:repro}

All code, certificates, result records and proof-object tooling are public
at \url{https://github.com/Leo-Y-Zhang/MathRecords}. The lower-bound
certificates check in milliseconds with no solver installed:

\begin{center}\small
\texttt{python vdw/verify\_certificate.py "<colouring>" <j> <t1> <t2>}
\end{center}

\noindent The encoding audit (\texttt{vdw/encoding\_audit.py}), the scale
test (\texttt{vdw/scale\_test.py}), the published-value replays
(\texttt{vdw/vdw\_validate.py}) and the refutations
(\texttt{vdw/vdw\_probe.py}) run with Python and PySAT; the repository's
\texttt{verify\_all.py} gate runs the full battery, including the DRAT
ladder when \texttt{kissat} and \texttt{drat-trim} binaries are present.
The headline proof-object certification is reproduced by
\texttt{vdw/cube\_certify.py} and \texttt{vdw/cube\_exhaustive.py} (about
four hours on four cores, resumable). Build notes for the two binaries,
including two silent Windows-specific pitfalls that produce wrong results
rather than errors, are in \texttt{vdw/DRAT.md}.

\begin{thebibliography}{99}

\bibitem{ahmed2009}
T.~Ahmed,
\emph{Some new van der Waerden numbers and some van der Waerden-type numbers},
Integers \textbf{9} (2009), A06, 65--76.

\bibitem{ahmed2011}
T.~Ahmed,
\emph{On computation of exact van der Waerden numbers},
Integers \textbf{11} (2011), \#A71.

\bibitem{ahmed2013}
T.~Ahmed,
\emph{Some more van der Waerden numbers},
Journal of Integer Sequences \textbf{16} (2013), Article 13.4.4.

\bibitem{aks}
T.~Ahmed, O.~Kullmann and H.~Snevily,
\emph{On the van der Waerden numbers $w(2;3,t)$},
Discrete Applied Mathematics \textbf{174} (2014), 27--51.

\bibitem{totalizer}
O.~Bailleux and Y.~Boufkhad,
\emph{Efficient CNF encoding of Boolean cardinality constraints},
in: Principles and Practice of Constraint Programming (CP 2003),
Lecture Notes in Computer Science \textbf{2833}, Springer, 2003, 108--122.

\bibitem{cadical}
A.~Biere, K.~Fazekas, M.~Fleury and M.~Heisinger,
\emph{CaDiCaL, Kissat, Paracooba, Plingeling and Treengeling entering the
SAT Competition 2020},
in: Proc.\ of SAT Competition 2020 -- Solver and Benchmark Descriptions,
Department of Computer Science Report Series B, vol.\ B-2020-1,
University of Helsinki, 2020, 50--53.

\bibitem{dransfield}
M.~R.~Dransfield, L.~Liu, V.~W.~Marek and M.~Truszczy\'nski,
\emph{Satisfiability and computing van der Waerden numbers},
The Electronic Journal of Combinatorics \textbf{11}(1) (2004), \#R41.

\bibitem{cubeandconquer}
M.~J.~H.~Heule, O.~Kullmann, S.~Wieringa and A.~Biere,
\emph{Cube and conquer: guiding CDCL SAT solvers by lookaheads},
in: Hardware and Software: Verification and Testing (HVC 2011),
Lecture Notes in Computer Science \textbf{7261}, Springer, 2012, 50--65.

\bibitem{pythagorean}
M.~J.~H.~Heule, O.~Kullmann and V.~W.~Marek,
\emph{Solving and verifying the Boolean Pythagorean triples problem via
cube-and-conquer},
in: Theory and Applications of Satisfiability Testing (SAT 2016),
Lecture Notes in Computer Science \textbf{9710}, Springer, 2016, 228--245.

\bibitem{pysat}
A.~Ignatiev, A.~Morgado and J.~Marques-Silva,
\emph{PySAT: a Python toolkit for prototyping with SAT oracles},
in: Theory and Applications of Satisfiability Testing (SAT 2018),
Lecture Notes in Computer Science \textbf{10929}, Springer, 2018, 428--437.

\bibitem{kourilpaul}
M.~Kouril and J.~L.~Paul,
\emph{The van der Waerden number $W(2,6)$ is 1132},
Experimental Mathematics \textbf{17} (2008), no.~1, 53--61.

\bibitem{landmanrobertson}
B.~M.~Landman and A.~Robertson,
\emph{Ramsey Theory on the Integers}, 2nd edition,
Student Mathematical Library \textbf{73}, American Mathematical Society,
2014.

\bibitem{oeis}
OEIS Foundation Inc.,
\emph{The On-Line Encyclopedia of Integer Sequences},
\url{https://oeis.org}, 2026. Sequences
\href{https://oeis.org/A217005}{A217005},
\href{https://oeis.org/A217007}{A217007},
\href{https://oeis.org/A217058}{A217058},
\href{https://oeis.org/A217059}{A217059},
\href{https://oeis.org/A217236}{A217236}.

\bibitem{vdw1927}
B.~L.~van der Waerden,
\emph{Beweis einer Baudetschen Vermutung},
Nieuw Archief voor Wiskunde \textbf{15} (1927), 212--216.

\bibitem{drattrim}
N.~Wetzler, M.~J.~H.~Heule and W.~A.~Hunt, Jr.,
\emph{DRAT-trim: efficient checking and trimming using expressive clausal
proofs},
in: Theory and Applications of Satisfiability Testing (SAT 2014),
Lecture Notes in Computer Science \textbf{8561}, Springer, 2014, 422--429.
Tool at \url{https://github.com/marijnheule/drat-trim}.

\end{thebibliography}

\end{document}
